% ===================================================================
%  جدول نمبر 9 — پہاڑ۔ جنگل۔
%
%  Traduction, faite en 2026, en ourdou d'aujourd'hui, sur l'ido de
%  texto/io. Regles de la colonne : voir 10-jadval-01.tex. Deux
%  scenes, deux numerotations : les cles portent c1 et c2. Ecrit avec
%  outils/ordo.py sous les yeux.
%
%  LA PLANCHE DE LA MONTAGNE EST CELLE OU L'OURDOU A LE MOINS BESOIN
%  D'EMPRUNTER, ET C'EST LA GEOGRAPHIE QUI L'EXPLIQUE. Le tamoul
%  avait du emprunter dix-sept mots ici, le telougou neuf, le coreen
%  quatre : la mesure des trois colonnes sur cette meme planche
%  formait une courbe monotone qu'on ne pouvait pas tracer avec moins
%  de trois colonnes. L'ourdou en emprunte QUATRE — مارموٹ la
%  marmotte (39), میگپائی la pie (56), شامو le chamois (54),
%  گلیشیئر le glacier (5) — et il rejoint le coreen tout en bas de
%  la courbe. La raison n'est pas la parente des langues : c'est que
%  le Pendjab regarde le Karakoram. Le col, le glacier, l'avalanche,
%  le torrent, le sapin, le cedre, l'aigle, le berger d'altitude sont
%  des choses ourdoues, et la langue les nomme depuis toujours :
%  درّہ, برفانی تودہ, صنوبر, دیودار, عقاب.
%
%  ET POURTANT گلیشیئر EST UN EMPRUNT, DANS LA LANGUE D'UN PAYS QUI
%  PORTE LE BALTORO ET LE SIACHEN. C'est le contre-exemple le plus
%  net du releve : avoir la chose sous les yeux ne suffit pas a lui
%  donner un nom propre. L'ourdou nomme le glacier par l'anglais des
%  cartes d'etat-major, parce que c'est l'alpinisme qui l'a decrit
%  avant l'ecole. Le mot suit le LIVRE ici, pas la carte — l'inverse
%  du marron chaud du tableau 7.
%
%  LES DEUX EMPRUNTS QUI RESTENT SONT DES BETES D'EUROPE. La
%  marmotte (39) et la pie (56) vivent bien dans l'Himalaya, mais
%  elles n'ont pas de nom qui soit descendu dans la plaine ; on les
%  nomme donc a l'anglaise. Le chamois (54), lui, n'existe pas du
%  tout en Asie : on aurait pu ecrire مارخور, qui est le grand
%  caprin du Pakistan et son animal national — c'est exactement le
%  piege du houblon et du vin. On ne l'a pas fait : la chose du
%  fac-simile est un chamois des Pyrenees. Huitieme refus de ce
%  genre.
%
%  LE PERSAN SAVANT PORTE TOUTE LA MINE ET TOUTE LA PIERRE :
%  رصدگاہ l'observatoire (12), سنگِ سماق le porphyre, بیروزہ la
%  resine (27), افعی la vipere, خرجی la besace (40), قلعہ le fort
%  (8), درّہ le defile (27). Et کوس rend la lieue — le kos vaut
%  environ trois kilometres, soit la lieue a peu de chose pres, la
%  ou le tamoul avait du prendre காதம், qui en vaut de dix a seize.
%  Deuxieme colonne a mesurer la lieue, deuxieme resultat, et c'est
%  l'ourdou qui tombe juste.
%
%  UN OISEAU QUI SE NOMME PAR SON TRAVAIL : کٹھ پھوڑا le pic (30),
%  « celui qui fend le bois ». Le tamoul avait relève la meme facon
%  de nommer pour la bergeronnette, « l'oiseau qui remue la queue ».
%  Deux langues sans lien, un meme procede, et les deux fois c'est
%  l'oiseau qu'on entend ou qu'on voit travailler qui garde un nom
%  vivant.
%
%  SIX FAUTES AU PREMIER JET, ET CINQ D'ENTRE ELLES SE REPARENT AVEC
%  LE MEME OUTIL. Le tableau 7 avait montre que l'ourdou possede la
%  relative postposee « جو ... ہے » et qu'elle evite de couper la
%  phrase en deux. Ici elle fait mieux que depanner : elle est la
%  reponse ordinaire chaque fois que l'ido pose le renvoi sur la
%  TETE avant le renvoi sur le complement — la scie (34) avant le
%  bucheron (33), la cognee (38) avant le journalier (39), le cor
%  (53) avant les cavaliers (52), le troupeau (50) avant le paturage
%  (49). En ourdou a tete finale, l'ordre naturel est l'inverse dans
%  les quatre cas ; la relative postposee le retablit sans rien
%  forcer. Une colonne a tete finale qui possede cet outil ne devrait
%  plus jamais rendre plus de deux divergences par tableau. On le
%  verifiera au tableau 10.
%
%  ET DEUX FAUTES QUI NE SONT PAS DE LANGUE, DEJA VUES AILLEURS :
%  une cle inventee — t09-c1-05-2, la ou l'ido continue le meme bloc
%  apres son saut de page — et un mot de brouillon laisse dans le
%  corps du fichier. La premiere avait ete faite au tableau 9 tamoul,
%  la seconde aux tableaux 5 et 7 tamouls. Les deux ont ete prises
%  par les controles, la cle par renvoji.py et le bloc ampute par
%  kolonoj.py (« 30% de la longueur mediane »), mais aucune n'aurait
%  du etre ecrite : on relit le bloc entier avant de passer au
%  suivant, et on remplace au lieu d'ajouter.
%
%  TRENTE-CINQ PAIRES, LE COMPTE LE PLUS HAUT DU VOLUME APRES LE 16.
%  Le tableau decrit plus qu'il ne raconte : le bloc c2-09-1 pose a
%  lui seul douze renvois, et le paysage de foret en aligne huit.
%  C'est la SYNTAXE qui coute, une fois de plus : ces blocs-la
%  attachent au lieu d'enumerer.
% ===================================================================

%%K t09-apar-1 apar
\VUsaut{4.0mm}
\VUtitre{15.0pt}{60}{جدول نمبر 9}
\VUsaut{3.0mm}

%%K t09-tit-1 sub
\VUcentre{12.0pt}{0}{\textit{پہلا منظر۔}}
\VUsaut{3.0mm}

%%K t09-tit-2 sub
\VUpk{11.4pt}{40}{پہاڑ۔ --- چشموں کا شہر۔}
\VUsaut{3.0mm}

%%K t09-c1-01-1 p
1. --- آٹھ دن سے میں پراٹس میں ہوں۔ جب میں پیرینیز کے حیرت انگیز
\VUgras{پہاڑی سلسلے} \textsuperscript{(1)} کے سامنے کھڑا ہوا تو جو
تحسین میں نے محسوس کی، اُسے پورا بیان نہیں کر سکتا ؛ اُس سلسلے کا
\VUgras{کنگرہ} \textsuperscript{(2)} نیلے آسمان پر کسی دیو قامت جھالر
کی طرح دکھائی دیتا ہے۔

%%K t09-c1-02-1 p
2. --- مجھے گمان بھی نہ تھا کہ پیرینیز کی \VUgras{چوٹیاں}
\textsuperscript{(3)} اتنی \VUgras{بلندی} تک پہنچتی ہیں، اور میں اُن
شاندار \VUgras{گلیشیئروں} \textsuperscript{(5)} اور اُن برف پوش
\VUgras{نوکیلی چوٹیوں} \textsuperscript{(4)} کے سامنے دنگ رہ گیا جن پر
سے اتنے ہولناک \VUgras{برفانی تودے} لڑھکتے ہیں۔

%%K t09-tit-3 sub
\VUsaut{3.0mm}
\VUpk{10.2pt}{40}{پہاڑی منظر۔}
\VUsaut{3.0mm}

%%K t09-c1-03-1 p
3. --- پہنچنے کے اگلے ہی دن میں اُس بجھے ہوئے پرانے
\VUgras{آتش فشاں} \textsuperscript{(6)} تک گیا جو پراٹس کے قریب ہے۔
\VUgras{دہانے} کے بنجر اور ننگے کناروں پر اُس \VUgras{لاوے} کے بہنے
کے نشان اب بھی صاف دکھائی دیتے ہیں جو کئی صدیاں پہلے \VUgras{پہاڑ کے
پہلو} \textsuperscript{(7)} پر بہا تھا۔ ایک \VUgras{ڈھلان} پر مجھے اُس
لاوے کے چند ٹکڑے مل بھی گئے۔

%%K t09-c1-04-1 p
4. --- پراٹس سرحد پر واقع ہے، اور جو \VUgras{قلعہ} \textsuperscript{(8)}
فرانس اور ہسپانیہ کو جدا کرنے والے \VUgras{درّے} \textsuperscript{(27)}
کی حفاظت کرتا ہے، وہ بالکل دریائے رائن کے کنارے کے اُن پرانے قلعوں کی
یاد دلاتا ہے جو اپنی چٹان کی چوٹی سے \VUgras{شکار} تاکتے لگتے ہیں۔

%%K t09-c1-05-1 p
5. --- ایک بہت بڑا \VUgras{عقاب} \textsuperscript{(9)}، جس کا گھونسلا
\VUgras{قلعے} \textsuperscript{(8)} سے دور نہیں، اکثر میری توجہ کھینچ
لیتا ہے۔

وہ \VUgras{شکاری پرندہ}، جس کی \VUgras{چونچ} مضبوط ہے، اکثر اپنے
بچوں کے لیے میمنا یا بکری کا بچہ اٹھا لاتا ہے اور اُسے اپنے
\VUgras{پنجوں} میں جکڑے رکھتا ہے۔ اُس کے \VUgras{پر} اتنے بڑے ہیں کہ
وہ اپنے بھاری بوجھ کے ساتھ دیر تک ہوا میں تیر سکتا ہے۔

%%K t09-tit-4 sub
\VUsaut{3.0mm}
\VUpk{10.2pt}{40}{سیاح۔}
\VUsaut{3.0mm}

%%K t09-c1-06-1 p
6. --- وہاں سب سے خوشگوار سیر « کاؤ » کے آبشار \textsuperscript{(10)}
تک کی ہے۔ \VUgras{پانی کا جھرنا} گھومتا ہوا گرتا ہے اور پھر شہر کے
مغرب میں واقع \VUgras{صنوبر کے جنگل} \textsuperscript{(11)} میں ایک
خوبصورت \VUgras{ندی} بن کر غائب ہو جاتا ہے۔ ہم اِسی جنگل میں سے
\VUgras{موسمیاتی رصدگاہ} \textsuperscript{(12)} تک چڑھے، اور
\VUgras{نظارہ گاہ} کی چوٹی سے ہمیں سارے علاقے کا \VUgras{پورا نظارہ}
دکھائی دیا۔ \VUgras{منظر} حیرت انگیز ہے، اور \VUgras{قدرت} اپنے سارے
خزانے اِس سرزمین پر لٹا دیتی لگتی ہے۔

%%K t09-c1-07-1 p
7. --- اترنے کے لیے ہم نے \VUgras{کیبل کار} \textsuperscript{(13)}
استعمال کی اور ایک چھوٹے \VUgras{اسٹیشن} پر رُکے، اُس پرانے
\VUgras{جاگیردار قلعے} کے \VUgras{کھنڈروں} \textsuperscript{(14)} کے
پاس جس میں کبھی پراٹس کے سردار رہتے تھے۔

%%K t09-c1-08-1 p
8. --- بے چارہ قلعہ! نیچے اُس بانکے \VUgras{چشموں کے شہر}
\textsuperscript{(15)} کو دیکھ کر وہ کیا سوچتا ہو گا، جو اب فیشن کا
\VUgras{گرم چشموں کا مرکز} بن چکا ہے؟

%%K t09-c1-08-2 p
\VUgras{آب و ہوا} اتنی خوشگوار ہے اور \VUgras{معدنی پانی} اتنا اثر
والا کہ \VUgras{شفا خانے} \textsuperscript{(17)} میں ہونے والا
\VUgras{علاج} سچ مچ ایک لطف ہے۔

%%K t09-tit-5 sub
\VUsaut{3.0mm}
\VUpk{10.2pt}{40}{خوشگوار چشمہ نگر۔}
\VUsaut{3.0mm}

%%K t09-c1-09-1 p
9. --- ویسے بھی قیام کو خوشگوار بنانے کے لیے سب کچھ کیا گیا ہے۔ ریل
ممکن بنانے کے لیے ایک بڑا \VUgras{اونچا پُل} \textsuperscript{(16)}
بنایا گیا ؛ \VUgras{حمام گاہ} کا \VUgras{باغ} \textsuperscript{(18)}،
جہاں \VUgras{محفلِ موسیقی} ہوتی ہے، بڑے دل کش انداز میں سجایا گیا
ہے۔ کئی خوش نما « \VUgras{کوٹھیاں} » \textsuperscript{(19)} جوا
خانے \textsuperscript{(21)} کے گرد بنی ہیں، اور ایک بہت اچھی
حالت میں رکھی \VUgras{پختہ سڑک} \textsuperscript{(22)} آمد و رفت کو
بہت آسان کر دیتی ہے۔

%%K t09-c1-10-1 p
10. --- ایک سادہ \VUgras{مٹی کا بند} \textsuperscript{(23)}
\VUgras{سڑک} \textsuperscript{(20)} کو اصل \VUgras{وادی}
\textsuperscript{(25)} سے جدا کرتا ہے ؛ اُس وادی کی تہ میں ایک خوش نما
چھوٹی \VUgras{جھیل} \textsuperscript{(26)} پڑی ہے، جو ایک صاف شفاف
\VUgras{ندی} \textsuperscript{(24)} کو پانی دیتی ہے۔

%%K t09-c1-11-1 p
11. --- وادی کی دوسری طرف \VUgras{لوہے کی کان} \textsuperscript{(28)}
سے کام لیا جاتا ہے۔ میں کئی بار \VUgras{کان کنوں} کو \VUgras{کان کے
کنویں} میں اترتے دیکھنے گیا، اور اُس \VUgras{سرنگ} میں بھی گھسا جس
میں وہ اپنا دن گزارتے ہیں، اُس \VUgras{کچی دھات} کو نکالتے ہوئے جس
میں \VUgras{دھات} ہوتی ہے۔

%%K t09-c1-12-1 p
12. --- کان کے قریب ایک بڑا \VUgras{ڈھلائی گھر} بنایا گیا ہے تاکہ
وہاں سے نکلنے والی دولت فوراً کام میں آئے۔ \VUgras{بڑی بھٹّی}
\textsuperscript{(29)}، جو یہاں کثرت سے ملنے والے \VUgras{پتھر کے
کوئلے} سے گرم کی جاتی ہے، \VUgras{ڈھلا لوہا} جلد اور سستا حاصل کرنے
کے قابل بناتی ہے۔

%%K t09-c1-13-1 p
13. --- کان سے دور نہیں، ایک \VUgras{پتھر کی کان}
\textsuperscript{(30)} کھودی جاتی ہے۔ بوڑھا \VUgras{پتھر کن} اپنی
\VUgras{کدال} سے نہ \VUgras{گرینائٹ} کاٹتا ہے، نہ \VUgras{سنگِ سماق}،
نہ \VUgras{ریتلا پتھر} ہی، بلکہ صرف گھر بنانے کے لیے سادہ
\VUgras{تراشا پتھر}۔

%%K t09-c1-14-1 p
14. --- اِس سرزمین کے سب سے خوبصورت زیوروں میں سے ایک
\VUgras{صنوبر} \textsuperscript{(31)} ہیں۔ ہر جگہ — \VUgras{کھڈ}
\textsuperscript{(32)} کی تہ میں، \VUgras{تیز نالے}
\textsuperscript{(33)} کے کنارے، \VUgras{اتھاہ گہرائی}
\textsuperscript{(34)} یا \VUgras{کھائی} کے کنارے — اُن کی باریک
سوئیاں سبز رنگ بکھیرتی ہیں۔

%%K t09-tit-6 sub
\VUsaut{3.0mm}
\VUpk{10.2pt}{40}{پیدل سیر۔}
\VUsaut{3.0mm}

%%K t09-c1-15-1 p
15. --- میں اپنی چھٹیوں سے فائدہ اٹھا کر \VUgras{لمبی} سیر کرتا ہوں۔
جو \VUgras{راستے} \textsuperscript{(35)} پہاڑ پر بل کھاتے ہیں، وہ
فاصلے کی پروا کیے بغیر کئی \VUgras{کلومیٹر} اور اکثر کئی
\VUgras{کوس} طے کرنے دیتے ہیں۔

%%K t09-c1-15-2 p
\VUgras{سنگِ میل} \textsuperscript{(36)} اور \VUgras{راہ نما کھمبے}
\textsuperscript{(37)} اُن راستوں کو نشان زد کرتے ہیں جن پر اکثر کوئی
جوان \VUgras{پہاڑی} \textsuperscript{(38)} ملتا ہے، پہلو میں
\VUgras{خرجی} \textsuperscript{(40)} لٹکائے اور \VUgras{مارموٹ}
\textsuperscript{(39)} اٹھائے ؛ یا کوئی \VUgras{بنجارا}
\textsuperscript{(41)}، جس کی \VUgras{کھال کی ٹوپی}
\textsuperscript{(42)} اُس کے پرانے پھٹے \VUgras{کُلاہ}
\textsuperscript{(43)} سے عجیب طرح میل نہیں کھاتی۔ وہ ایک سدھایا ہوا
\VUgras{زنجیر بندھا ریچھ} \textsuperscript{(44)} ہانکتا ہے۔

%%K t09-tit-7 sub
\VUpk{10.2pt}{40}{کوہ پیمائی۔}
\VUsaut{3.0mm}

%%K t09-c1-16-1 p
16. --- تقریباً روز میں \VUgras{رہبر} \textsuperscript{(48)} کے ساتھ
\VUgras{چڑھائی} کی مشق کرنے جاتا ہوں، اور مجھے بڑا فخر ہوتا ہے جب
پیٹھ پر \VUgras{پِٹھو} \textsuperscript{(47)} باندھے اور ہاتھ میں «
\VUgras{الپن اسٹاک} » لیے، کسی سچے \VUgras{سیاح} \textsuperscript{(46)}
کی طرح میں \VUgras{چٹانوں} \textsuperscript{(45)} پر حملہ آور ہوتا
ہوں۔

%%K t09-c1-17-1 p
17. --- پہاڑ کی چوٹی سے میدان کے باشندے چیونٹیوں سے بڑے نہیں لگتے ؛
\VUgras{بھیڑوں کا ریوڑ} \textsuperscript{(50)}، جو \VUgras{چراگاہ}
\textsuperscript{(49)} میں چر رہا ہے، بچوں کا کھلونا لگتا ہے۔ کوئی
\VUgras{چیڑ} \textsuperscript{(51)} یا کوئی \VUgras{لکڑی کا گھروندا}
\textsuperscript{(52)} ہریالی پر بمشکل ایک دھبے سا دکھائی دیتا ہے۔

%%K t09-c1-18-1 p
18. --- اِن سیروں کے دوران ہمیں اکثر \VUgras{پگڈنڈی}
\textsuperscript{(53)} پر کوئی \VUgras{شامو کا شکاری}
\textsuperscript{(54)} ملتا ہے، کندھے پر وہ جانور اٹھائے جسے اُس نے
ابھی مارا ہے۔

%%K t09-c1-19-1 p
19. --- میں نے اپنے پودوں کے دفتر میں \VUgras{سرخس}
\textsuperscript{(55)} کی ایک نہایت قابلِ ذکر شاخ رکھی ہے، اور
\VUgras{ہیدر} \textsuperscript{(57)} کی ایک خوبصورت گلابی ٹہنی، جو
میں نے اپنی پچھلی سیر میں ایک چھوٹے \VUgras{غار} \textsuperscript{(56)}
کے دہانے پر توڑی تھیں۔

%%K t09-tit-8 sub
\VUsaut{3.0mm}
\VUcentre{12.0pt}{0}{\textit{دوسرا منظر۔}}
\VUsaut{3.0mm}

%%K t09-tit-9 sub
\VUpk{11.4pt}{40}{جنگل۔ --- شکار۔}
\VUsaut{3.0mm}

%%K t09-c2-01-1 p
1. --- کل میں نے جنگل میں ایک لذیذ دن گزارا۔ جن جانوروں اور
\VUgras{کیڑوں} \textsuperscript{(5)} سے وہ آباد ہے، اُن میں جو
مستعدی چھائی ہوئی ہے، اُس نے مجھے بہت دل چسپ لگا۔

%%K t09-c2-01-2 p
\VUgras{مکڑی} \textsuperscript{(1)}، جو دو شاخوں کے بیچ اپنا
\VUgras{جالا} \textsuperscript{(2)} بُنتی ہے تاکہ گزرتی
\VUgras{مکھی} \textsuperscript{(3)} پکڑ سکے ؛ \VUgras{گھونگا}
\textsuperscript{(4)}، جو بڑی مشکل سے اپنا خول ڈھوتا ہے ؛ سینگوں والا
بھونرا، جو اپنے شکار کا جائزہ لے رہا ہے ؛ مستعد چیونٹی، جو
\VUgras{چیونٹیوں کے بھٹ} \textsuperscript{(6)} کے پاس بڑے جوش میں ہے —
یہ سب ننھے جاندار غیر معمولی مستعدی سے آتے، جاتے اور کام کرتے رہے۔

%%K t09-c2-02-1 p
2. --- ایک \VUgras{سانپ} \textsuperscript{(7)}، جسے میں نے پہلی نظر
میں \VUgras{افعی} سمجھا مگر جو سادہ \VUgras{دھامن} کے سوا کچھ نہ
تھا، درخت کے تنے کے نیچے دبکا ہوا تھا اور وقفے وقفے سے اپنا باریک سا
سر باہر نکالتا تھا، جو ہر بار ڈسنے کو تیار لگتا تھا۔ میرے سامنے ایک
موٹا \VUgras{بے خول گھونگا} \textsuperscript{(8)} بھاری چال سے رینگ
رہا تھا، اُن مزے دار \VUgras{جنگلی اسٹرابیریوں} \textsuperscript{(11)}
میں سے ایک نگل کر جو وہاں کثرت سے اُگتی ہیں۔

%%K t09-c2-03-1 p
3. --- زمین \VUgras{جنگلی پھولوں} سے بچھی ہوئی تھی ؛
\VUgras{پرائم روز} \textsuperscript{(9)}، \VUgras{سدا بہار}
\textsuperscript{(10)} اور بہت سے دوسرے پودے، جنھیں میں نہیں جانتا
تھا، \VUgras{گھاس کے گچھوں} \textsuperscript{(12)} کے بیچ کھلے ہوئے
تھے۔

%%K t09-c2-04-1 p
4. --- اُس علاقے میں \VUgras{ٹرفل} تقریباً اتنی ہی عام ہے جتنی
\VUgras{کھمبی} \textsuperscript{(14)}، اور ایک جنگل بان کا
\VUgras{جنگلی سؤر کا بچہ} \textsuperscript{(13)} چٹورے پن سے ٹٹول رہا
تھا کہ شاید اُسے کچھ مل جائے۔

%%K t09-tit-10 sub
\VUsaut{3.0mm}
\VUpk{10.2pt}{40}{چوری کا شکار۔}
\VUsaut{3.0mm}

%%K t09-c2-05-1 p
5. --- میں ایک ایسے منظر کے \VUgras{اختتام} پر پہنچا جس نے مجھے بہت
محظوظ کیا۔ اپنے \VUgras{باسٹ کتّے} \textsuperscript{(15)} کی سونگھ کی
مدد سے \VUgras{جنگل بان} \textsuperscript{(16)} نے ابھی ابھی ایک
\VUgras{چور شکاری} \textsuperscript{(22)} کو عین اُس وقت جا لیا تھا جب
وہ اپنا مارا ہوا \VUgras{شکار} \textsuperscript{(17)} اٹھا کر لے جانے
کی تیاری کر رہا تھا۔ گرفتار ہونے سے بہتر اپنا مال چھوڑ دینا سمجھ کر
چور شکاری بھاگ نکلا تھا، اور ایک \VUgras{جنگلی خرگوش}
\textsuperscript{(18)}، ایک \VUgras{ہرنی} \textsuperscript{(19)}، ایک
\VUgras{کاکڑ} \textsuperscript{(20)} اور ایک \VUgras{جنگلی سؤر}
\textsuperscript{(21)} گھاس پر پڑے جنگل بان کو خوش کر رہے تھے۔

%%K t09-c2-06-1 p
6. --- جنگل میں جتنی قسم کے درخت ہیں، وہ حیران کر دیتی ہے۔
\VUgras{بیچ کے درخت} \textsuperscript{(23)} اور \VUgras{بھوج پتر}
\textsuperscript{(26)}، \VUgras{چیڑ}، جو \VUgras{بیروزہ}
\textsuperscript{(27)} دیتے ہیں، \VUgras{بلوط} \textsuperscript{(29)}
اور اِن کے سوا بہت سے اور — سب اپنی شاخیں ایک دوسرے میں الجھائے ہوئے
ہیں۔

%%K t09-c2-06-2 p
\VUgras{آئیوی} \textsuperscript{(24)}، جو اونچے درختوں کے تنے سے چمٹی
رہتی ہے، \VUgras{بن} \textsuperscript{(25)} کو چمکیلی ہریالی سے رنگ
دیتی ہے ؛ وہ ہریالی \VUgras{کائی} \textsuperscript{(31)} کے ہلکے اور
پھیکے سبز رنگ سے خوشگوار طور پر مل جاتی ہے، اور اُسی کائی میں
\VUgras{کٹھ پھوڑا} \textsuperscript{(30)} کیڑے تلاش کرتا ہے۔

%%K t09-tit-11 sub
\VUsaut{3.0mm}
\VUpk{10.2pt}{40}{جنگل کا منظر۔}
\VUsaut{3.0mm}

%%K t09-c2-07-1 p
7. --- پرانے بلوطوں نے مجھے موہ لیا۔ اُن کی موٹی بل کھائی
\VUgras{جڑیں} \textsuperscript{(32)}، جو زمین سے آدھی نکلی ہوئی ہیں،
اُنھیں طاقت اور توانائی کا شان دار روپ دیتی ہیں، اور میں خود سے
پوچھتا ہوں کہ \VUgras{مالک} \textsuperscript{(35)} اُنھیں کٹوا کر اُس
\VUgras{آری} \textsuperscript{(34)} کے حوالے کرنے پر کیسے راضی ہو
جاتا ہے جو \VUgras{لکڑہارے} \textsuperscript{(33)} کی ہے۔ وہ بے
چارے \VUgras{تنے} \textsuperscript{(36)}، جو اپنی \VUgras{چھال}
\textsuperscript{(37)} سے محروم کر دیے گئے یا اُس
\VUgras{کلہاڑی} \textsuperscript{(38)} سے چیر دیے گئے جو
\VUgras{دیہاڑی دار مزدور} \textsuperscript{(39)} کی ہے، مجھے دُکھ
دیتے ہیں۔

%%K t09-c2-08-1 p
8. --- پاس ہی ایک بوڑھے \VUgras{کوئلہ پکانے والے}
\textsuperscript{(41)} نے اپنے \VUgras{بیلچے} سے کوئلے کا ایک
\VUgras{ڈھیر} \textsuperscript{(42)} تیار کیا تھا، اور ایک بڑھیا کٹے
ہوئے درختوں کی ٹہنیوں کی \VUgras{سوکھی لکڑی} کا \VUgras{گٹھا}
\textsuperscript{(40)} اٹھا لائی تھی۔

%%K t09-tit-12 sub
\VUsaut{3.0mm}
\VUpk{10.2pt}{40}{شکار۔}
\VUsaut{3.0mm}

%%K t09-c2-09-1 p
9. --- میں چند لمحے آرام کرنے \VUgras{جنگل بان کے گھر}
\textsuperscript{(46)} جانا چاہتا تھا ؛ مگر جب میں اُس \VUgras{چھوٹی
جھیل} \textsuperscript{(43)} کے پاس پہنچا جس پر ایک خوبصورت
\VUgras{ہنس} \textsuperscript{(45)} تیرتا ہے، تو میں نے دیکھا کہ ایک
حالیہ \VUgras{سیلاب} \textsuperscript{(44)} نے راستے کو سچ مچ کی
\VUgras{دلدل} میں بدل دیا ہے۔ مجھے مڑنا پڑا، اور جنگل کے کنارے کھڑے
\VUgras{دیودار} \textsuperscript{(49)}، \VUgras{سفیدوں}
\textsuperscript{(48)} اور \VUgras{نارون} \textsuperscript{(47)} کے
پاس سے گزرتے ہوئے میں نے دیکھا کہ ایک \VUgras{جھاڑی زار}
\textsuperscript{(54)} سے ایک نہایت خوبصورت \VUgras{بارہ سنگھا}
\textsuperscript{(50)} نکلا، جس کے پیچھے ایک ضدی \VUgras{شکاری کتّوں
کا غول} \textsuperscript{(51)} لگا ہوا تھا ؛ اُس غول کو وہ
\VUgras{بگل} \textsuperscript{(53)} بھی بھڑکا رہا تھا جو
\VUgras{گھڑ سوار شکاری} \textsuperscript{(52)} بجا رہے تھے۔ بے چارا جانور بالکل نڈھال
تھا۔ وہ آ کر میرے چند قدم کے فاصلے پر مرتے مرتے گر پڑا۔

%%K t09-c2-10-1 p
10. --- جنگل بان کا بیٹا، جسے \VUgras{دوڑے شکار} کے شور نے کھینچ لیا
تھا، گھر سے نکلا اور ہم دونوں پھر جنگل میں چلے گئے۔ اُس نے ایک پرانے
بلوط کی شاخ پر ایک \VUgras{میگپائی} \textsuperscript{(56)} دیکھی تھی۔
اُس کا خیال تھا کہ اُس کا گھونسلا شاید \VUgras{پتوں کے جھنڈ}
\textsuperscript{(58)} میں چھپا ہے، اور وہ گھونسلا پکڑنا چاہتا تھا۔

%%K t09-c2-10-2 p
\VUgras{گلہری} \textsuperscript{(61)} کی طرح پھرتیلا، وہ \VUgras{شاخ}
\textsuperscript{(55)} سے شاخ پر چڑھتا گیا، مگر اُسے کچھ نہ ملا اور وہ
صرف ایک بوڑھے \VUgras{کوّے} \textsuperscript{(60)} کو اُڑانے میں کامیاب
ہوا، جو ایک \VUgras{ٹہنی} \textsuperscript{(57)} پر بیٹھا تھا۔

\VUsaut{4.0mm}
\VUfilet{22mm}
